% foundations/system_layers.tex

\chapter{System Layers}

\label{chap:system-layers}

% ============================================================
\section*{Motivation}
% ============================================================

\begin{remark}
Railway alignment modelling combines multiple fundamentally different
domains.
\end{remark}

\begin{remark}
Classical approaches often merge:
\begin{itemize}
\item intrinsic geometry,
\item coordinate reference systems,
\item metric interpretation,
\item construction realization,
\item runtime evaluation,
\item dynamics,
\item and operational vehicle interpretation
\end{itemize}
into a single implicit model.
\end{remark}

\begin{remark}
Within AIM, these concerns are separated into interacting conceptual and
operational system layers.
\end{remark}

\begin{remark}
This separation improves:
\begin{itemize}
\item conceptual clarity,
\item numerical robustness,
\item scalability,
\item interoperability,
\item runtime flexibility,
\item and realization-aware modelling.
\end{itemize}
\end{remark}

% ============================================================
\section{Overview}
% ============================================================

\begin{definition}[System layer]
\label{def:system-layer}
A system layer is a coherent subsystem operating on its own:
\begin{itemize}
\item state representation,
\item operators,
\item constraints,
\item interpretation rules,
\item and runtime assumptions.
\end{itemize}
\end{definition}

\begin{remark}
The AIM framework separates railway alignment modelling into:
\begin{itemize}
\item geometry layer,
\item CRS layer,
\item metric realization layer,
\item physical realization layer,
\item runtime layer,
\item vehicle layer,
\item and optimization layer.
\end{itemize}
\end{remark}

\begin{remark}
These layers are conceptually distinct, but operationally coupled.
\end{remark}

% ============================================================
\section{Geometry Layer}
% ============================================================

\subsection{Role}

\begin{remark}
The geometry layer defines the intrinsic semantic structure of the
alignment.
\end{remark}

\begin{remark}
Its central quantities are:
\begin{itemize}
\item arc-length \(s\),
\item curvature \(\kappa(s)\),
\item cant evolution,
\item transition structure,
\item and pose evolution laws.
\end{itemize}
\end{remark}

% ------------------------------------------------------------

\subsection{Core Relations}

\begin{remark}
The geometry layer is governed by the canonical intrinsic relations
defined in \cref{chap:core-equations}.
\end{remark}

\begin{remark}
In particular, it uses:
\begin{itemize}
\item the curvature relation \cref{eq:curvature-relation},
\item the reconstruction relation \cref{eq:state-evolution},
\item and the cant evolution relation \cref{eq:cant-evolution}.
\end{itemize}
\end{remark}

% ------------------------------------------------------------

\subsection{Characteristics}

\begin{remark}
The geometry layer is:
\begin{itemize}
\item intrinsic,
\item semantic,
\item continuous,
\item analytical,
\item and independent of a particular metric backend.
\end{itemize}
\end{remark}

\begin{remark}
It defines alignment semantics rather than physically realized track
geometry.
\end{remark}

% ============================================================
\section{CRS Layer}
% ============================================================

\subsection{Role}

\begin{remark}
The CRS layer defines external coordinate reference systems and their
relations to the Earth.
\end{remark}

\begin{remark}
Typical coordinate systems include:
\begin{itemize}
\item engineering CRS,
\item projected CRS,
\item geodetic coordinates,
\item local construction coordinate systems,
\item and ellipsoid-based representations.
\end{itemize}
\end{remark}

% ------------------------------------------------------------

\subsection{Coordinate Transformation}

\begin{remark}
CRS transformations are represented by the coordinate transformation
operator \(\mathcal C\), introduced in
\cref{def:coordinate-transform-operator}.
\end{remark}

\begin{remark}
The CRS layer operates within world space:
\[
\mathcal C
:
\Omega_{\mathrm{world}}
\rightarrow
\Omega_{\mathrm{world}}.
\]
\end{remark}

% ------------------------------------------------------------

\subsection{Fundamental Observation}

\begin{proposition}
\label{prop:alignment-not-crs-invariant}
A railway alignment is generally not invariant under CRS
transformations.
\end{proposition}

\begin{remark}
In particular:
\begin{itemize}
\item straight lines,
\item curvature,
\item transition behaviour,
\item stationing-related length interpretation,
\item and cant-related spatial interpretation
\end{itemize}
may change under projection or CRS-dependent embedding.
\end{remark}

\begin{remark}
The CRS layer therefore affects spatial interpretation, but it is not
identical to intrinsic railway projection.
\end{remark}

% ============================================================
\section{Metric Realization Layer}
% ============================================================

\subsection{Role}

\begin{remark}
The metric realization layer defines how semantic alignment structures
are interpreted geometrically within a selected metric model.
\end{remark}

\begin{remark}
It is governed conceptually by the metric realization operator
\(\mathcal G\), introduced in
\cref{def:metric-realization-operator}.
\end{remark}

% ------------------------------------------------------------

\subsection{Core Principle}

\begin{principle}[Metric realization principle]
\label{princ:metric-realization-layer}
Semantic alignment logic is separated from metric realization behaviour.
\end{principle}

\begin{remark}
Thus:
\begin{itemize}
\item semantic curvature laws,
\item transition structure,
\item sparse alignment logic,
\item and stationing semantics
\end{itemize}
remain conceptually stable even when metric realization changes.
\end{remark}

% ------------------------------------------------------------

\subsection{Metric Variants}

\begin{remark}
Typical metric realization modes include:
\begin{itemize}
\item engineering-Euclidean realization,
\item projection-aware realization,
\item ellipsoid-aware realization,
\item and future differential-geometric realization.
\end{itemize}
\end{remark}

% ------------------------------------------------------------

\subsection{Operator Boundary}

\begin{remark}
The metric realization layer forms the operator boundary between:
\begin{itemize}
\item semantic alignment logic,
\item and metric-dependent geometric evaluation.
\end{itemize}
\end{remark}

\begin{remark}
This architecture is introduced formally in
\cref{chap:metric-realization}.
\end{remark}

% ============================================================
\section{Physical Realization Layer}
% ============================================================

\subsection{Role}

\begin{remark}
The physical realization layer models the transformation from analytical
target geometry to physically realized railway geometry.
\end{remark}

% ------------------------------------------------------------

\subsection{Core Operator}

\begin{remark}
The physical realization layer is governed by the realization operator
\[
\mathcal R
:
\mathcal X_{\mathrm{target}}
\rightarrow
\mathcal X_{\mathrm{real}},
\]
introduced in \cref{def:realization-operator}.
\end{remark}

% ------------------------------------------------------------

\subsection{Characteristics}

\begin{remark}
The physical realization layer includes:
\begin{itemize}
\item construction constraints,
\item machine behaviour,
\item discretization,
\item interpolation,
\item structural response,
\item maintenance processes,
\item and measurement feedback.
\end{itemize}
\end{remark}

\begin{remark}
Physical realization acts approximately as a low-pass filter on geometry
and curvature.
\end{remark}

\begin{remark}
Within AIM, physical realization is interpreted as an explicit operator
layer rather than as a secondary engineering detail.
\end{remark}

\begin{remark}
This layer is developed in
\cref{chap:construction-and-realization}.
\end{remark}

% ============================================================
\section{Runtime Layer}
% ============================================================

\subsection{Role}

\begin{remark}
The runtime layer transforms stored railway representations into
operationally evaluable geometry.
\end{remark}

\begin{remark}
Runtime evaluation is governed by the evaluation operator
\[
\mathcal E
:
\mathcal X
\times
\Omega_{\mathrm{track}}
\rightarrow
\mathcal Y,
\]
introduced in \cref{def:evaluation-operator}.
\end{remark}

% ------------------------------------------------------------

\subsection{Runtime Services}

\begin{remark}
Typical runtime services include:
\begin{itemize}
\item world--track projection,
\item frame evaluation,
\item metric-aware geometry evaluation,
\item visualization,
\item BIM placement,
\item spatial queries,
\item sensor integration,
\item and dynamic quantity evaluation.
\end{itemize}
\end{remark}

% ------------------------------------------------------------

\subsection{Runtime Characteristics}

\begin{remark}
The runtime layer is:
\begin{itemize}
\item state-dependent,
\item scale-dependent,
\item metric-aware,
\item operator-based,
\item computationally adaptive,
\item and runtime-evaluated.
\end{itemize}
\end{remark}

\begin{remark}
Many runtime quantities are evaluated dynamically rather than stored
explicitly.
\end{remark}

\begin{remark}
The runtime layer therefore acts as the operational execution layer of
the AIM framework.
\end{remark}

\begin{remark}
The runtime architecture is developed further in
\cref{chap:runtime-services}.
\end{remark}

% ============================================================
\section{Vehicle Layer}
% ============================================================

\subsection{Role}

\begin{remark}
The physically relevant object is not merely the alignment centerline,
but the railway vehicle moving within space.
\end{remark}

% ------------------------------------------------------------

\subsection{Operational Geometry}

\begin{remark}
The vehicle layer includes:
\begin{itemize}
\item vehicle envelopes,
\item bogie motion,
\item wheelset orientation,
\item center-of-gravity motion,
\item operational clearances,
\item and platform interaction.
\end{itemize}
\end{remark}

% ------------------------------------------------------------

\subsection{Runtime Dependency}

\begin{remark}
Vehicle-space interpretation depends on runtime-evaluated pose3 states,
frames, cant, profile, and metric context.
\end{remark}

\begin{remark}
Thus, the vehicle layer is not merely an application layer, but an
operational interpretation layer built on top of runtime services.
\end{remark}

% ------------------------------------------------------------

\subsection{Motivation}

\begin{remark}
Classical railway alignment systems often approximate operational
geometry using simplified clearance assumptions.
\end{remark}

\begin{remark}
Within AIM, operational vehicle-space geometry may instead be modelled
explicitly.
\end{remark}

% ============================================================
\section{Optimization Layer}
% ============================================================

\subsection{Role}

\begin{remark}
The optimization layer searches for railway states satisfying:
\begin{itemize}
\item geometric,
\item metric,
\item dynamic,
\item operational,
\item realization,
\item and regulatory constraints.
\end{itemize}
\end{remark}

% ------------------------------------------------------------

\subsection{Realization-Aware Optimization}

\begin{remark}
Within AIM, optimization may be evaluated on realized states:
\[
J(\mathcal R(x(u))).
\]
\end{remark}

\begin{remark}
This relation is part of the canonical optimization structure introduced
in \cref{eq:general-optimization}.
\end{remark}

% ------------------------------------------------------------

\subsection{Hybrid Nature}

\begin{remark}
Optimization combines:
\begin{itemize}
\item analytical models,
\item sparse representations,
\item discrete constraints,
\item runtime evaluation,
\item realization operators,
\item metric-dependent geometry evaluation,
\item and operational vehicle-space constraints.
\end{itemize}
\end{remark}

% ============================================================
\section{Layer Interactions}
% ============================================================

\begin{remark}
The system layers interact continuously.
\end{remark}

\begin{remark}
For example:
\begin{itemize}
\item intrinsic geometry feeds metric realization,
\item CRS influences metric interpretation and world embedding,
\item metric realization affects runtime geometry,
\item physical realization modifies target states into realized states,
\item runtime supports projection, visualization, and optimization,
\item optimization modifies target geometry,
\item vehicle-space interpretation depends on pose3, frames, and runtime,
\item and physical measurements should generally reference realized geometry.
\end{itemize}
\end{remark}

% ============================================================
\section{Canonical Transformation Structure}
% ============================================================

\begin{definition}[Canonical transformation structure]
\label{def:canonical-transformation-structure}
A simplified conceptual AIM structure is:
\[
\text{intrinsic geometry}
\rightarrow
\text{metric realization}
\rightarrow
\text{runtime evaluation}
\rightarrow
\text{vehicle-space interpretation}.
\]
\end{definition}

\begin{remark}
Physical realization is a separate transformation:
\[
\text{target state}
\rightarrow
\text{realized state}.
\]
\end{remark}

\begin{remark}
CRS embedding influences metric realization and world-space
representation, but it is not identical to intrinsic world--track
projection.
\end{remark}

\begin{remark}
Optimization closes the loop by modifying target geometry based on
runtime-evaluated and realization-aware objectives.
\end{remark}

% ============================================================
\section{Scale Dependence}
% ============================================================

\begin{remark}
Different layers dominate at different spatial and operational scales.
\end{remark}

\begin{remark}
For example:
\begin{itemize}
\item CRS effects dominate at large spatial scales,
\item metric realization dominates at projection transitions,
\item physical realization effects dominate at construction scales,
\item runtime costs dominate in large operational systems,
\item and vehicle-space effects dominate near platforms, switches, and
clearance-critical situations.
\end{itemize}
\end{remark}

% ============================================================
\section{Canonical Layer Separation}
% ============================================================

\begin{principle}[Canonical separation principle]
\label{princ:canonical-layer-separation}
Intrinsic geometry, CRS embedding, metric realization, physical
realization, runtime interpretation, operational behaviour, and
optimization should not be merged into a single implicit model.
\end{principle}

\begin{remark}
Many inconsistencies in classical railway systems arise from implicit
mixing of fundamentally different conceptual layers.
\end{remark}

% ============================================================
\section{Key Insight}
% ============================================================

\begin{proposition}
\label{prop:alignment-multi-layer-system}
Railway alignment modelling is not a single geometric problem, but a
multi-layer interaction system.
\end{proposition}

\begin{proposition}
\label{prop:operational-behaviour-emerges}
The operational behaviour of railway systems emerges from interactions
between:
\begin{itemize}
\item intrinsic geometry,
\item CRS embedding,
\item metric realization,
\item physical realization,
\item runtime evaluation,
\item vehicle-space interpretation,
\item and optimization.
\end{itemize}
\end{proposition}

% ============================================================
\section{Connection to AIM}
% ============================================================

\begin{summary}
The AIM framework separates railway alignment modelling into interacting
system layers.

This separation enables:
\begin{itemize}
\item rigorous intrinsic geometry,
\item CRS-aware embedding,
\item metric-aware realization,
\item physical realization-aware modelling,
\item scalable runtime systems,
\item explicit vehicle-space interpretation,
\item and realization-aware optimization.
\end{itemize}

Together, these layers form the conceptual architecture underlying the
AIM framework.
\end{summary}
